% GENERATED FILE. Edit canonical lessons, then run npm run generate:books.
% canonical-source-hash: fnv1a64:72b311d543647165
% canonical-lessons: ZH-C61-chang, ZH-W61-chang-long, ZH-C61-duan, ZH-C61-gao, ZH-C61-ai, ZH-C61-xin

\chapter{Long and Short}
\label{ch:zh-long-and-short}
\hlchaptermodality{61}

\begin{chapteropening}
\textbf{By the end of this chapter:} \emph{I can say long, short, tall, short of height, and new.}

Complete the last lesson of chapter 61: I can say long, short, tall, short of height, and new.
\end{chapteropening}

\section[\texorpdfstring{\zh{长} (cháng)}{长 (cháng)}]{\zh{长} (cháng) — long}

\label{lesson:ZH-C61-chang}



\begin{quote}
Before the new one: say the Mandarin for black, then the Mandarin for blue.
\end{quote}

\subsection*{You'll want to know: \zh{长}}
\textbf{\zh{长}} — \emph{cháng} — \textquotedblleft{}long\textquotedblright{}.

\begin{quote}
\emph{jīntiān de diànyǐng hěn cháng} — Today's film is a long one.
\end{quote}

The new character, in pieces: \zh{长} — four strokes: a short fall, a long bar, a vertical that flicks up at the foot, and a sweep down to the right. The next lesson puts it in your hand.

\begin{cousinweb}
\textbf{\zh{长}} (\emph{cháng}) is inherited from Middle Chinese \emph{drjang} (Guangyun, in Baxter's transcription), glossed 'lasting', 'far' and 'for ever'.
\end{cousinweb}

\begin{grammarlens}[title={What you've built}]
One more word for everyday answers.
\end{grammarlens}

\subsection*{Your turn}
\begin{itemize}
\raggedright
  \item \emph{Say it:} \emph{cháng}
  \item \emph{Say it:} \emph{cháng}, once more
  \item \emph{Say it:} \emph{jīntiān de diànyǐng hěn cháng} — Today's film is a long one.
  \item \emph{Recall:} say the Mandarin for black, then the Mandarin for blue, then say \emph{cháng} again
\end{itemize}

\begin{tcolorbox}[breakable,colback=teal!4,colframe=teal!35!black,title={Before you move on}]
What does \zh{长} mean? (\textquotedblleft{}Long\textquotedblright{}.) Say it once more.
\end{tcolorbox}
\section[\texorpdfstring{\zh{长} (cháng)}{长 (cháng)}]{\zh{长} — four strokes in one piece}

\label{lesson:ZH-W61-chang-long}



\begin{quote}
The word first, then the character inside it.

\begin{itemize}
\raggedright
  \item \emph{Recall:} say \emph{cháng} — \textbf{R1}, one lesson back, before it decays
\end{itemize}

\begin{quote}
\zh{大} \zh{八}
\end{quote}

\emph{Read it:} each one aloud, with nothing said first — \emph{dà}, \emph{bā}
\end{quote}

\subsection*{Script — four strokes in one piece}
\hlblockfigure{\detokenize{figures/ZH-W61-chang-long-filmstrip.pdf}}{How \zh{长} is written, stroke by stroke}

\begin{quote}
\zh{长}
\end{quote}

\textbf{One piece, four strokes.}

\textbf{A short left-falling stroke} at the top right, falling toward the left.

\textbf{A long horizontal} across the middle.

\textbf{A vertical} from the top left, down through the bar, that turns and flicks up to the right at the foot.

\textbf{A long stroke falling to the right}, from the bar down to the bottom right.

\subsection*{Your turn}
\begin{itemize}
\raggedright
  \item \emph{Write it:} \zh{长}, all four, with the model in front of you
  \item \emph{Cover:} write it again from memory, and count the pieces, not the strokes
\end{itemize}

\begin{tcolorbox}[breakable,colback=teal!4,colframe=teal!35!black,title={Before you move on}]
How many strokes in \zh{长}? (\textbf{Four}.) Which stroke flicks up at the foot? (\textbf{The vertical, the third}.) Which way does the last one fall? (\textbf{Down to the right}.)

Source: \href{https://raw.githubusercontent.com/chanind/hanzi-writer-data/68d10a4b21150cae5e1ebbd223eed289cf32d90c/data/\%E9\%95\%BF.json}{Hanzi Writer Data \zh{长}, pinned stroke paths and medians}.
\end{tcolorbox}
\section[duǎn]{duǎn — short in length; brief}

\label{lesson:ZH-C61-duan}



\begin{quote}
Before the new one: say the Mandarin for blue, then the Mandarin for long.
\end{quote}

\subsection*{You'll want to know: duǎn}
\emph{duǎn} — \textquotedblleft{}short in length; brief\textquotedblright{}.

\begin{quote}
\emph{wǒ de kùzi hěn duǎn} — My trousers are short.
\end{quote}

Its characters come later; for now it is a word for the ear.

\begin{cousinweb}
\emph{duǎn} is inherited from Middle Chinese \emph{twanX} (Guangyun, in Baxter's transcription), glossed 'hurried' and 'not long'.
\end{cousinweb}

\begin{grammarlens}[title={What you've built}]
One more word for everyday answers.
\end{grammarlens}

\subsection*{Your turn}
\begin{itemize}
\raggedright
  \item \emph{Say it:} \emph{duǎn}
  \item \emph{Say it:} \emph{duǎn}, once more
  \item \emph{Say it:} \emph{wǒ de kùzi hěn duǎn} — My trousers are short.
  \item \emph{Recall:} say the Mandarin for blue, then the Mandarin for long, then say \emph{duǎn} again
\end{itemize}

\begin{tcolorbox}[breakable,colback=teal!4,colframe=teal!35!black,title={Before you move on}]
What does duǎn mean? (\textquotedblleft{}Short in length; brief\textquotedblright{}.) Say it once more.
\end{tcolorbox}
\section[gāo]{gāo — tall; high}

\label{lesson:ZH-C61-gao}



\begin{quote}
Before the new one: say the Mandarin for long, then the Mandarin for short in length.
\end{quote}

\subsection*{You'll want to know: gāo}
\emph{gāo} — \textquotedblleft{}tall; high\textquotedblright{}.

\begin{quote}
\emph{wǒ gēge hěn gāo tā xǐhuan yùndòng} — My brother is tall, and he loves sport.
\end{quote}

Its characters come later; for now it is a word for the ear.

\begin{cousinweb}
\emph{gāo} is inherited from Middle Chinese \emph{kaw} (Guangyun, in Baxter's transcription), glossed 'up', 'lofty' and 'far'.
\end{cousinweb}

\begin{grammarlens}[title={What you've built}]
One more word for everyday answers.
\end{grammarlens}

\subsection*{Your turn}
\begin{itemize}
\raggedright
  \item \emph{Say it:} \emph{gāo}
  \item \emph{Say it:} \emph{gāo}, once more
  \item \emph{Say it:} \emph{wǒ gēge hěn gāo tā xǐhuan yùndòng} — My brother is tall, and he loves sport.
  \item \emph{Recall:} say the Mandarin for long, then the Mandarin for short in length, then say \emph{gāo} again
\end{itemize}

\begin{tcolorbox}[breakable,colback=teal!4,colframe=teal!35!black,title={Before you move on}]
What does gāo mean? (\textquotedblleft{}Tall; high\textquotedblright{}.) Say it once more.
\end{tcolorbox}
\section[ǎi]{ǎi — short in height; low}

\label{lesson:ZH-C61-ai}



\begin{quote}
Before the new one: say the Mandarin for short in length, then the Mandarin for tall.
\end{quote}

\subsection*{You'll want to know: ǎi}
\emph{ǎi} — \textquotedblleft{}short in height; low\textquotedblright{}.

\begin{quote}
\emph{wǒ mèimei hěn ǎi} — My little sister is short.
\end{quote}

Its characters come later; for now it is a word for the ear.

\begin{cousinweb}
\emph{ǎi} is inherited from Middle Chinese \emph{'eaX} (Guangyun, in Baxter's transcription), glossed 'short to look at'.
\end{cousinweb}

\begin{grammarlens}[title={What you've built}]
One more word for everyday answers.
\end{grammarlens}

\subsection*{Your turn}
\begin{itemize}
\raggedright
  \item \emph{Say it:} \emph{ǎi}
  \item \emph{Say it:} \emph{ǎi}, once more
  \item \emph{Say it:} \emph{wǒ mèimei hěn ǎi} — My little sister is short.
  \item \emph{Recall:} say the Mandarin for short in length, then the Mandarin for tall, then say \emph{ǎi} again
\end{itemize}

\begin{tcolorbox}[breakable,colback=teal!4,colframe=teal!35!black,title={Before you move on}]
What does ǎi mean? (\textquotedblleft{}Short in height; low\textquotedblright{}.) Say it once more.
\end{tcolorbox}
\section[xīn]{xīn — new}

\label{lesson:ZH-C61-xin}



\begin{quote}
Before the new one: say the Mandarin for tall, then the Mandarin for short in height.
\end{quote}

\subsection*{You'll want to know: xīn}
\emph{xīn} — \textquotedblleft{}new\textquotedblright{}.

\begin{quote}
\emph{zhè shì wǒ de xīn shǒujī} — This is my new phone.
\end{quote}

Its characters come later; for now it is a word for the ear.

\begin{cousinweb}
\emph{xīn} is inherited from Middle Chinese \emph{sin} (Guangyun, in Baxter's transcription), glossed 'new, as against old'.
\end{cousinweb}

\begin{grammarlens}[title={What you've built}]
One more word for everyday answers.
\end{grammarlens}

\subsection*{Your turn}
\begin{itemize}
\raggedright
  \item \emph{Say it:} \emph{xīn}
  \item \emph{Say it:} \emph{xīn}, once more
  \item \emph{Say it:} \emph{zhè shì wǒ de xīn shǒujī} — This is my new phone.
  \item \emph{Recall:} say the Mandarin for tall, then the Mandarin for short in height, then say \emph{xīn} again
\end{itemize}

\begin{tcolorbox}[breakable,colback=teal!4,colframe=teal!35!black,title={Before you move on}]
What does xīn mean? (\textquotedblleft{}New\textquotedblright{}.) Say it once more.
\end{tcolorbox}
